% =============================================================
% บทที่ 1 - บทนำ
% โครงงาน: การออกแบบและพัฒนาเกมเดินเรือด้วยระบบสั่งการด้วยเสียงเพื่อส่งเสริมทักษะการสื่อสาร
%          (Design and Development of a Voice-Controlled Ship Navigation Game to Enhance Communication Skills)
% =============================================================

% -----------------------------------------------------------
\section{ความเป็นมาและความสำคัญของปัญหา}
% -----------------------------------------------------------
ในปัจจุบัน ภาษาอังกฤษมีบทบาทสำคัญอย่างยิ่งในการเป็นเครื่องมือสื่อสารระดับสากล ทว่าการเรียนการสอนภาษาอังกฤษในห้องเรียนโดยทั่วไปยังคงเผชิญกับข้อจำกัดในการสร้างสภาพแวดล้อมที่เอื้อต่อการนำไปใช้งานจริง~\cite{sadigzade2025} ผู้เรียนจำนวนมากประสบปัญหาความวิตกกังวลในการใช้ภาษาต่างประเทศ (Foreign Language Anxiety) และกลัวความผิดพลาดเมื่อต้องเปล่งเสียงเพื่อสื่อสาร ส่งผลให้ขาดความมั่นใจและไม่สามารถเชื่อมโยงทักษะทางภาษาเข้ากับสถานการณ์ในชีวิตจริงหรือการโต้ตอบแบบทันท่วงทีได้~\cite{varma2025}
ที่ผ่านมา มีความพยายามนำสื่อดิจิทัลและเกมมาประยุกต์ใช้เพื่อการเรียนรู้ภาษา (Game-Based Language Learning) แต่ยังคงพบกับความท้าทายสำคัญ~\cite{sadigzade2025} งานวิจัยชี้ให้เห็นว่าแอปพลิเคชันหรือเกมเพื่อการเรียนรู้ภาษาแบบดั้งเดิมมักเน้นไวยากรณ์และการท่องจำคำศัพท์เป็นหลัก ทำให้รูปแบบการเล่นจำเจอยู่กับการคลิกเลือกคำตอบหรือการอ่านข้อความ ในขณะที่เกมที่เน้นความสนุกสนานก็มักไม่ได้กำหนดให้ผู้เล่นต้องเปล่งเสียงพูดออกมา~\cite{varma2025} ช่องว่างนี้แสดงให้เห็นถึงโอกาสในการสร้างสรรค์เกมที่สามารถบูรณาการการฝึกฝนทักษะการพูดเข้ากับบริบทที่น่าตื่นเต้นและท้าทาย เพื่อกระตุ้นให้เกิดการสื่อสารอย่างเป็นธรรมชาติ~\cite{zargham2024}
โครงงานการออกแบบและพัฒนาเกมเดินเรือด้วยระบบสั่งการด้วยเสียงเพื่อส่งเสริมทักษะการสื่อสารจึงจัดทำขึ้น เพื่อพัฒนาแนวทางการเรียนรู้ภาษาอังกฤษผ่านสภาพแวดล้อมแบบดื่มด่ำ (Immersive Environment) ผู้เล่นจะรับบทเป็น ``กัปตันเรือ'' ที่ต้องเผชิญกับสถานการณ์จำลองการเดินเรือฝ่าสภาพแวดล้อมที่ท้าทาย โดยภารกิจของผู้เล่นไม่ใช่เพียงการควบคุมทิศทางผ่านหน้าจอ แต่คือการใช้ ``เสียง'' เป็นเครื่องมือหลักในการสั่งการลูกเรือปัญญาประดิษฐ์ (AI Crew) ในตำแหน่งต่าง ๆ ได้แก่ ต้นกลเรือ (Chief Engineer) ผู้ควบคุมเครื่องยนต์และความเร็ว และพลถือท้ายและยามระวังหน้า (Helmsman and Lookout) ผู้บังคับหางเสือและเฝ้าระวังสิ่งกีดขวาง เพื่อบังคับเรือไม่ให้ชนสิ่งกีดขวางและนำเรือไปสู่เป้าหมายอย่างปลอดภัย โดยคำสั่งที่ใช้ในเกมยึดตามวลีสื่อสารมาตรฐานทางทะเลขององค์การทางทะเลระหว่างประเทศ (IMO Standard Marine Communication Phrases: SMCP)~\cite{imo2002smcp}

% -----------------------------------------------------------
\section{นิยามศัพท์เฉพาะ}
% -----------------------------------------------------------

\begin{enumerate}[label=\arabic*)]
    \item \textbf{เกมเพื่อการเรียนรู้ (Serious Game)} หมายถึง เกมที่ออกแบบโดยมีวัตถุประสงค์หลักเพื่อการเรียนรู้หรือการฝึกทักษะมากกว่าความบันเทิง
    \item \textbf{การรู้จำเสียง (Speech-to-Text / Speech Recognition)} หมายถึง เทคโนโลยีการแปลงคำสั่งเสียงพูดภาษาอังกฤษของผู้เล่นให้เป็นข้อความ เพื่อนำไปตรวจความถูกต้องและประมวลผลเป็นคำสั่งควบคุมเรือ
    \item \textbf{ตัวกรองทางอารมณ์ (Affective Filter)} หมายถึง อุปสรรคทางอารมณ์ เช่น ความกลัว ความวิตกกังวล และการขาดความมั่นใจ ซึ่งขัดขวางการเรียนรู้และการพูดภาษาที่สอง
    \item \textbf{ความวิตกกังวลในการใช้ภาษาต่างประเทศ} หมายถึง ความรู้สึกประหม่า กลัวการพูดผิด หรือกลัวการถูกประเมินในทางลบเมื่อต้องสื่อสารด้วยภาษาต่างประเทศ ในโครงงานนี้วัดด้วยแบบวัดความวิตกกังวลในการพูดภาษาอังกฤษที่ดัดแปลงจาก FLCAS
    \item \textbf{ลูกเรือปัญญาประดิษฐ์} หมายถึง ตัวละคร NPC (Non-Player Character) ในเกมที่รับคำสั่งเสียงจากผู้เล่นและให้เสียงตอบกลับ (Voice Feedback) เช่น การทวนคำสั่งหรือขอให้ย้ำคำสั่ง
    \item \textbf{สภาพแวดล้อมแบบดื่มด่ำ} หมายถึง สภาพแวดล้อมจำลองภายในเกมที่สร้างประสบการณ์การเรียนรู้ที่มีทั้งความท้าทาย ความสมจริง และความปลอดภัยทางความรู้สึก
\end{enumerate}

% -----------------------------------------------------------
\section{วัตถุประสงค์}
% -----------------------------------------------------------

\begin{enumerate}[label=\arabic*)]
    \item เพื่อออกแบบและพัฒนาเกมเดินเรือด้วยระบบสั่งการด้วยเสียง (Voice-Controlled Ship Navigation Game) สำหรับใช้เป็นสื่อการเรียนรู้แบบดื่มด่ำ (Immersive Learning) ที่มีความสมจริงและออกแบบเพื่อช่วยลดความวิตกกังวลในการใช้ภาษาต่างประเทศของผู้เรียน
    \item เพื่อประเมินประสิทธิภาพทางเทคนิคของระบบจากแฟ้มบันทึกผลในเกม ได้แก่ ความถูกต้องของการรู้จำคำสั่ง ร้อยละของการสื่อสารที่ลูกเรือเข้าใจ จำนวนครั้งที่ลูกเรือขอให้พูดซ้ำ และเวลาตอบสนอง และประเมินประสบการณ์การเล่นเกมของผู้เล่นด้วยแบบสอบถาม
    \item เพื่อเปรียบเทียบระดับความวิตกกังวลในการพูดภาษาอังกฤษของผู้เล่นก่อนและหลังการเล่นเกม โดยใช้แบบวัดที่ดัดแปลงจาก Foreign Language Classroom Anxiety Scale (FLCAS)
\end{enumerate}

% -----------------------------------------------------------
\section{เป้าหมายและขอบเขต}
% -----------------------------------------------------------

\subsection{ขอบเขตด้านเนื้อหา}

ระบบที่พัฒนาเป็นเกมแนวจำลองสถานการณ์ (Simulation Game) แบบผู้เล่นคนเดียว (Single-player) ที่ใช้เสียงเป็นช่องทางควบคุมหลัก ประกอบด้วยระบบหลัก ได้แก่

\begin{enumerate}[label=\arabic*)]
    \item \textbf{ระบบรู้จำเสียงและประมวลผลคำสั่ง (Speech Recognition \& Command Parser System):} ระบบรับเสียงจากไมโครโฟน แปลงเป็นข้อความ แล้วส่งให้ตัวแยกคำสั่ง (Command Parser) ตรวจสอบและแยกคำสั่งภาษาอังกฤษ
    \item \textbf{ระบบฟิสิกส์และการควบคุมเรือ (Ship Navigation Physics):} ระบบจำลองการเคลื่อนที่ของเรือที่มีความเฉื่อยและการตอบสนองต่อหางเสือที่สมจริง ทำให้ผู้เล่นต้องวางแผนล่วงหน้าในการออกคำสั่ง
    \item \textbf{ระบบลูกเรือปัญญาประดิษฐ์ (AI Crew System):} ระบบตัวละคร NPC ที่รับคำสั่งเสียงจากผู้เล่นและให้เสียงตอบกลับ เช่น การทวนคำสั่งเมื่อรับทราบ หรือการขอให้พูดซ้ำเมื่อไม่เข้าใจคำสั่ง
    \item \textbf{ระบบสร้างสิ่งกีดขวาง (Obstacle Generation System):} ระบบสร้างสิ่งกีดขวาง ได้แก่ หิน แนวหินโสโครก สันดอน และทุ่นระเบิด แบบสุ่มในเส้นทางการเดินเรือ ยกเว้นภารกิจโหมดเนื้อเรื่องซึ่งใช้แผนที่คงที่ เพื่อท้าทายให้ผู้เล่นต้องตัดสินใจออกคำสั่งหลบหลีกในทันที
    \item \textbf{ระบบส่วนติดต่อผู้ใช้ (UI/HUD):} หน้าจอแสดงข้อมูลที่จำเป็น เช่น เรดาร์ตรวจจับสิ่งกีดขวาง เข็มทิศแสดงองศาเรือ มาตรวัดความเร็ว และสัญลักษณ์แสดงสถานะการรับเสียงของไมโครโฟน
    \item \textbf{ระบบช่วยสอนภาษา (Language Support):} คำแนะนำประโยคที่ควรพูดในแต่ละสถานการณ์ การทวนคำสั่งของลูกเรือด้วยรูปประโยคมาตรฐาน และรายงานภาษาในหน้าสรุปผลเมื่อจบการเดินเรือ
    \item \textbf{โหมดเนื้อเรื่อง (Story Mode):} ภารกิจตามเนื้อเรื่อง 3 บท รวม 9 ภารกิจ
    \item \textbf{ระบบบันทึกผลการทดลอง (Research Log):} การบันทึกข้อมูลการเล่นลงแฟ้ม CSV แยกตามรอบการเปิดเกม เพื่อใช้วิเคราะห์ผล
\end{enumerate}

\subsection{ขอบเขตด้านประชากรและกลุ่มตัวอย่าง}

ประชากรคือนักศึกษาระดับปริญญาตรี กลุ่มตัวอย่างคืออาสาสมัครที่เลือกแบบสะดวก (Convenience Sampling) จำนวน 10--20 คน

\subsection{ขอบเขตด้านเทคนิค}

\begin{enumerate}[label=\arabic*)]
    \item พัฒนาตัวเกมด้วยเอนจิน Unity (Unity 6) และใช้ภาษา C\# สำหรับการพัฒนาสภาพแวดล้อมเกมและระบบฟิสิกส์ของการเดินเรือ
    \item นำเสนอภาพในรูปแบบมุมมองการเล่นแบบไอโซเมตริก (Isometric view) และรองรับการทำงานบนระบบปฏิบัติการ Windows PC
    \item ใช้ไมโครโฟนรับเสียงพูดภาษาอังกฤษเป็นช่องทางหลักในการสั่งการลูกเรือ โดยมีการพิมพ์คำสั่งผ่านแป้นพิมพ์เป็นช่องทางสำรอง ทั้งสองช่องทางผ่านขั้นตอนการประมวลผลคำสั่งเดียวกัน และไม่มีปุ่มบังคับหางเสือหรือคันโทรเลขเครื่องยนต์โดยตรง
    \item ใช้โมเดลการรู้จำเสียง Whisper ที่ทำงานภายในเครื่องของผู้เล่น (On-device) ผ่านไลบรารี whisper.cpp เพื่อแปลงเสียงพูดเป็นข้อความด้วยโมเดล ggml-base.en
    \item ใช้ API ของ OpenAI (โมเดล GPT-4o mini) เฉพาะเมื่อตัวแยกคำสั่งไม่เข้าใจคำสั่ง และเมื่อผู้เล่นสนทนากับลูกเรือ
\end{enumerate}

% -----------------------------------------------------------
\section{ประโยชน์ที่คาดว่าจะได้รับ}
% -----------------------------------------------------------

\begin{enumerate}[label=\arabic*)]
    \item ผู้เล่นได้ฝึกใช้วลีสั่งการมาตรฐานทางทะเลในสถานการณ์จำลอง
    \item ผู้เล่นมีความมั่นใจและกล้าพูดภาษาอังกฤษมากขึ้น
    \item ผู้สอนสามารถนำเกมไปใช้เป็นสื่อเสริมการเรียนการสอนเชิงปฏิบัติ
\end{enumerate}